\begin{abstract}
We construct a new family, indexed by the odd integers $N\geq 1$, of
 $(2+1)$-dimensional quantum field theories called {\it quantum
 hyperbolic field theories} (QHFT), and we study its main structural
 properties. The QHFT are defined for (marked) $(2+1)$-bordisms
 supported by compact oriented $3$-manifolds $Y$ with a properly
 embedded framed tangle $L_\Ff$ and an {\it arbitrary}
 $PSL(2,\C)$-character $\rho$ of $Y \setminus L_\Ff$ (covering, for
 example, the case of hyperbolic cone manifolds). The marking of QHFT
 bordisms includes a specific set of parameters for the space of
 pleated hyperbolic structures on punctured surfaces.  Each QHFT
 associates in a constructive way to any triple $(Y,L_\Ff,\rho)$ with
 marked boundary components a tensor built on the matrix dilogarithms,
 which is holomorphic in the boundary parameters. We establish {\it
 surgery formulas} for QHFT partitions functions and describe their relations
 with the {\it quantum hyperbolic invariants} of \cite{BB1,BB2}
 (either defined for unframed links in closed manifolds and characters
 trivial at the link meridians, or hyperbolic {\it cusped}
 3-manifolds). For every $PSL(2,\mc)$-character of a punctured
 surface, we produce new families of conjugacy classes of ``moderately
 projective" representations of the mapping class groups.
\end{abstract}

\smallskip

\emph{{\bf Keywords:} hyperbolic geometry, quantum field theory;
mapping class groups, quantum invariants, Cheeger-Chern-Simons
class, dilogarithms.}

\section{Introduction}\label{INTRO}

In this paper we construct a new family $\{\Hh_N\}$, indexed by the odd
integers $N\geq 1$, of $(2+1)$-dimensional quantum field theories
(QFT) that we call Quantum Hyperbolic Field Theories (QHFT). Here,
following \cite{At,T}, by  QFT we mean a functor from a
$(2+1)$-bordism category, possibly non purely topological, to the
tensorial category of finite dimensional complex linear spaces.
\smallskip

The QHFT bordism category is based on triples $(Y,L_\Ff,\rho)$, where
$Y$ is a compact oriented $3$-manifold, possibly with non-empty
boundary $\partial Y$, $L_\Ff$ is a properly embedded {\it framed
tangle} (ie. a framed $1$-dimensional non oriented submanifold) in
$Y$, and $\rho$ is a flat $sl(2,\mc)$-connection on $Y \setminus
L_\Ff$ up to gauge equivalence (ie. a $PSL(2,\C)$-character of $Y
\setminus L_\Ff$), with {\it arbitrary} holonomy at the meridians of
the tangle components. We require furthermore that $L_\Ff$ is
non-empty when $N>1$, and that it intersects each of the boundary
component, if any.  

We will also consider a variant, denoted QHFT$^0$, such that the
tangles $L$ are {\it unframed}, while the characters $\rho$ are
defined on the whole of $Y$, that is the meridian holonomies are
trivial. Finally we consider a ``fusion'' of QHFT and QHFT$^0$
(still denoted QHFT) that incorporates both, by considering tangles
having a {\it framed} part $L_\Ff$ as well as an {\it unframed} one
$L^0$ (see Section \ref{variant}).

The objects of the bordism category are suitably {\it marked}
surfaces. Every such a  QHFT {\it surface} is a diffeomorphism $f:
(S,\Tt,p(\beta)) \to \Sigma$, where $\Tt$ is a so called
``efficient triangulation" of a fixed base oriented surface $S$
with genus $g$ and $r$ marked framed points $p_i$ ($r>0$ and $r>2$
if $g=0$), and $\beta$ is any $PSL(2,\C)$-character of $S\setminus
\{p_i\}$, represented by points $p(\beta)$ in specific parameter
spaces for ${\rm Hom}(\pi,PSL(2,\C))$, built on $\Tt$ and
particularly suited to the QHFT. In fact we construct several such
parameter spaces with small ``residual gauge groups'' acting on
them, and we point out the relations to each other. One of them,
the so called $(-)$-{\it exponential $\Ii$-parameter space}, is
defined in terms of cross-ratios and incorporates the
Bonahon-Thurston shearbend coordinates for pleated hyperbolic
surfaces with punctures.

Every QHFT bordism has marked boundary and is considered as a
``transition'' from its input QHFT surfaces towards the output
ones. We understand that the characters $\rho$ and $\beta$ are
compatible. Every QHFT functor associates to such a transition a
tensor called the {\it amplitude}, defined up to a sign and
multiplication by $N$th roots of unity.
\smallskip

When $Y=W$ is closed (that is $\partial Y = \emptyset$), the
amplitudes $\Hh_N(W,L_\Ff, \rho)$ are numerical invariants called {\it
  partition functions}. The QHFT$^0$ partition functions
$H_N(W,L,\rho)$ coincide with the ``quantum hyperbolic invariants''
constructed in \cite{BB1, BB2}, while the QHFT ones yield new wide
families of numerical invariants, covering interesting geometric
situations, such as compact hyperbolic cone manifolds. We will analyze
the relations between $\Hh_N$ and $H_N$ partition functions.  In
\cite{BB2} we defined also quantum hyperbolic invariants $H_N(M)$ for
non-compact complete hyperbolic 3-manifolds $M$ of finite volume,
ie. for {\it cusped manifolds}. Although these invariants are not
immediately QHFT partition functions, we will show how they can be
obtained in terms of these last. For that we establish a ``surgery
formula'' for quantum hyperbolic invariants of cusped manifolds and
QHFT partition functions that generalizes the one for
Cheeger-Chern-Simons classes, and makes a crucial use of some of
W. Neumann's arguments in \cite{N}, sections 11 and 14.
\smallskip

By restricting QHFT to the trivial bordisms (the cylinders) we get
a new family of conjugacy classes of ``moderately projective''
representations of the mapping class groups of punctured surfaces,
that is, defined up to a sign and multiplication by $N$th roots of
unity.
\smallskip

We stress that we need that any bordism includes a non-empty link,
intersecting each boundary component (so that the QHFT surfaces
have punctures), in order to build a consistent functor when $N>1$
(see \cite{BB2}, Lemma 6.4). Even when the holonomy is trivial
around the punctures, we cannot forget them, in particular for
what concerns the mapping class groups.
\smallskip

We show that QHFT are in fact restrictions to a {\it geometric bordism
  category} of ``universal functors'' called Quantum Hyperbolic
Geometry (QHG). QHG includes the definition of a specific category of
triangulated 3-dimensional pseudomanifolds equipped with additional
structures, and modeled on the functional properties of the {\it
  matrix dilogarithms} studied in \cite{BB2}.  The QHG functors
associate determined tensors to every such a decorated triangulation,
obtained by tracing the matrix dilogarithms supported by each
tetrahedron. The key point is that such tensors are invariant up to
QHG triangulated pseudomanifold isomorphism.  The main step in order
to construct specializations with a strong geometric content, such as
QHFT, consists in converting each QHFT (marked) bordism to a QHG
triangulated pseudomanifold, unique up to QHG triangulated
pseudomanifold isomorphism.
\smallskip

Hence we view this paper as a kind of achievement of the fundation of
the theory initiated in \cite{BB1, BB2}.
\smallskip

A main interest in the QHFT comes from the fact that they relate
classical $3$-dimensional hyperbolic geometry to the world of quantum
field theories, two main themes of low-dimensional topology that
remained essentially disjoint since their spectacular developments in
the early eighties. In particular, the celebrated Kashaev's Volume
Conjecture for hyperbolic knots in $S^3$ \cite{K2} appears as a
special instance of the challenging general problem of understanding
the relations between the asymptotic behaviour of QHFT partition
functions and fundamental invariants coming from differential
geometry, like the Cheeger-Chern-Simons class (see \cite{BB1}, section
5, \cite{BB2}, section 7, and section \ref{ALL} of the present
paper). We plan to face the asymptotics of QHFT partition functions in
future works.
\smallskip

 In \cite{B} the spectrum of the mapping class group representations
 of Section \ref{MODREP} is studied by using geometric quantization of
 the Bonahon-Thurston complex intersection $2$-form. 

A natural problem left unsettled is to determine the relations between
the QHFT and Turaev's Homotopic QFT \cite{Tu2}. This and formulas
describing the behaviour of the QHFT amplitudes under framing changes
will be treated in a sequel to this paper, as they rely mainly on
$R$-matrix computations.

We refer to \cite{BBo} for a discussion about QHG in the framework of
gravity in dimension $3$.

\begin{remark}\label{OKCCS}
{\rm The results of this paper can be repeated almost verbatim to
define Cheeger-Chern-Simons invariants for QHFT bordisms, by replacing
the matrix dilogarithms for $N>1$ with Neumann's extended Rogers
dilogarithm, corresponding to $N=1$ (see Remark \ref{scissors}). This
is described with all details in \cite{BB2} for the case of cusped
manifolds and triples $(W,L,\rho)$ (ie. for quantum hyperbolic
invariants). Hence in the following we concentrate on the quantum
theory $N>1$, which is technically harder.}
\end{remark}

Here is the content of the paper.
\smallskip

The universal QHG functors are defined in section \ref{QHG}, where
we recall also from \cite{BB1,BB2} and \cite{N0,N} (section
\ref{ToGeoSpec} in particular) the notions and results we need. 
\smallskip

The QHFT bordism category is described in section \ref{BORDCAT}, while
its objects, the QHFT surfaces, are developed starting with
section \ref{PARAM}.
\smallskip

The QHFT functors are defined in section \ref{TENSOR}. This includes
the construction of the distinguished QHG triangulated pseudomanifolds
associated to any triple $(Y,L_\Ff,\rho)$ with marked boundary
components, and of the trace tensors computed on them.
The conjugacy classes of moderately projective
representations of the mapping class groups are treated in section
\ref{MODREP}.
\smallskip

The partition functions $\Hh_N(W,L_\Ff,\rho)$ are considered in
section \ref{CLOSED}.  We show in section \ref{ALL} that when $\rho$
is defined on the whole of $W$, $\Hh_N(W,L_\Ff,\rho)$ coincides with
$H_N(W,L\cup L',\rho)$, where $L'$ is a parallel copy of the unframed
link $L$ given by the framing $\Ff$. In section \ref{DEHNF} we prove
the surgery formula for quantum hyperbolic invariants of cusped
manifolds and QHFT partition functions. In fact the QHG
pseudomanifolds used to compute the trace tensors carry certain
{\it cohomological weights} (see Section \ref{ToGeoSpec}) and the 
partition function values actually depend also on them. These weights
play indeed a subtle role in the surgery formulas.  This eventually
leads to realize $H_N(M)$ as the limit of $H_N(M_n,L_n,\rho_n)$, where
$M_n$ is a sequence of closed hyperbolic manifolds converging
geometrically to $M$, $L_n$ is the link of geodesic cores of the
hyperbolic Dehn fillings of $M$ that produce $M_n$, and $\rho_n$ is
the hyperbolic holonomy of $M_n$.  In section \ref{FIBERED} we discuss
alternative computations of the QHFT partition functions for manifolds
that fiber over $S^1$.  For fibred cusped manifolds, this allows in
particular to identify each $H_N(M)$ with a special instance of QHFT
partition function. We conclude with an example in section
\ref{EXAMPLES}.
\medskip

{\bf Convention.}  Unless otherwise stated all manifolds are oriented,
and the boundary is oriented via the convention: {\it last is the
  ingoing normal}. Often we denote ``$=_N$'' the equality of tensors
up to sign and multiplication by $N$th roots of unity.

\section{Universal QHG}\label{QHG}

\subsection{Building blocks}\label{buildbox} The building blocks of 
QHG are {\it flat/charged $\Ii$-tetrahedra} $(\Delta,b,w,f,c)$, and
{\it matrix dilogarithms} $\Rr_N(\Delta(b,w,f,c)) \in {\rm Aut}(\mc^N
\otimes \mc^N)$ defined for every odd positive integer $N$ \cite{BB2}.


\medskip

\noindent{\bf Flat/charged-$\Ii$ tetrahedra.} Consider the half-space
model of the oriented hyperbolic space $\mh^3$, with the group of
direct isometries identified with $PSL(2,\C)$ by the conformal action
on $\partial \overline {\mh}^3= \mc\mathbb{P}^1 = \C \cup\{ \infty \}$
by Moebius transformations. An $\Ii$-{\it tetrahedron} is an oriented
ideal tetrahedron $\Delta$ in $\partial \overline{\mh}^3$ with distinct
ordered vertices $v_0$, $v_1$, $v_2$ and $v_3$ on $\partial
\overline{\mh}^3$.

In fact we consider $\Delta$ as an abstract oriented simplex
equipped with an additional decoration. The ordering of the
vertices is encoded by a {\it branching} $b$, that is, edge
orientations obtained via the rule: each edge points towards the
biggest end-point. Each $2$-face has an induced branching, and a
{\it $b$-orientation}, which is just compatible with that of two
edges on the boundary. We order the $2$-faces $\delta_0$, $\ldots
\delta_3$ by the opposite vertices, and the edges $e_0$, $e_1$,
$e_2$ of $\delta_3$ by stipulating that for $j=0,1$, $v_j$ is the
first end-point of $e_j$.  For exactly two $2$-faces the
$b$-orientation and the boundary orientation are the same. The
{\it $b$-orientation} of $\Delta$ coincides with the given one if
the $b$-orientation of $\delta_3$ looks anti-clockwise from $v_3$.
We give $\Delta$ and each $2$-face $\delta$ a $b$-{\it sign} $*_b$
and $\sigma (\delta)$ respectively, which is $1$ if the two
orientations agree, and $-1$ otherwise.

The hyperbolic structure is encoded by the cross-ratio moduli that
label the edges of $\Delta$.  Recall that opposite edges share the
same cross-ratio moduli. We set $w=(w_0,w_1,w_2)$ with
$w_j=w(e_j) \in \C\setminus \{0,1\}$.
Hence $w_{j+1} = 1/(1-w_j)$ (indices mod($\mz/3\mz$), and $$w_0 = (
v_2- v_1)(v_3-v_0)/(v_2-v_0)(v_3- v_1).$$ We say that the
$\Ii$-tetrahedron $(\Delta,b,w)$ is {\it non-degenerate} if it is of
non zero volume, that is if the imaginary part of each $w_i$ is not
zero; then they have the same sign $*_w = \pm 1$.
\smallskip

It is very convenient to encode $(\Delta,b,w)$ in dual terms. In
Figure \ref{CQDidealtensor} we show the $1$-skeleton of the dual
cell decomposition of Int$(\Delta)$ ($x$ and the indices $i$, $j$,
$k$ and $l$ are considered below). It is understood that (dual) edges
without arrows are incoming at the crossing. Note that an oriented
edge is outgoing exactly when the $b$-sign of the dual $2$-face is
$1$.
\begin{figure}[ht]
\begin{center}
 \includegraphics[width=8cm]{CQDidealtensor.eps}
\caption{\label{CQDidealtensor} $\Ii$-tetrahedra and dual
encoding.}
\end{center}
\end{figure}

A {\it flat/charged} $\Ii$-tetrahedron is an $\Ii$-tetrahedron
equipped with a {\it flattening} $f$ and a {\it charge} $c$, two
notions first introduced in \cite{N0} and \cite{N}. Flattenings
and charges are $\Z$-valued functions defined on the
edges of $\Delta$ that take the same value on opposite edges and
satisfy the following properties, respectively (where $\log$ has
the imaginary part in $]-\pi,\pi]$):
\begin{enumerate}
\item[F.] {\it Flattening condition:} ${\rm l}_0 + {\rm l}_1 +{\rm
    l}_2 = 0$, where
\begin{equation}\label{clb}
{\rm l}_j = {\rm l}_j(b,w,f)= \log(w_j) + \sqrt{-1}\pi f_j,
\end{equation}
\item[C.] {\it Charge condition:} $c_0+c_1+c_2=1$.
\end{enumerate}
We call ${\rm l}_j$ a {\it (classical) log-branch}, and for every
odd $N>1$ we define the \emph{(level $N$) quantum log-branch} as
\begin{equation}\label{qlb}
{\rm l}_{j,N} = \log(w_j) + \sqrt{-1}\pi (N+1)(f_j -*_bc_j).
\end{equation}
The bijective map $$({\rm l}_0,{\rm l}_1,{\rm l}_2) \mapsto
\left(w_0;\frac{{\rm l}_0-\log(w_0)}{\sqrt{-1}\pi},\frac{{\rm
l}_1-\log(w_1)}{\sqrt{-1}\pi}\right)$$ yields an identification of the
set of log-branches on $(\Delta,b)$ with the Riemann surface
$\widehat{\mc}$ of the maps $w_0\mapsto
(\log(w_0)+\varepsilon\pi\sqrt{-1},\log((1-w_0)^{-1})+
\varepsilon'\pi\sqrt{-1})$, with $\varepsilon,\varepsilon'\in
\{0,1\}$. Similarly, the set of triples $(w_0',w_1',w_2')$ with
$w_j'=\exp({\rm l}_{j,N}/N)$ gets identified with the quotient
covering by $N\mz\times N\mz$. These two spaces are the domains of
definition for the matrix dilogarithms to be described next, when
$N=1$ and $N>1$ respectively.
\medskip

\noindent {\bf Matrix dilogarithms.}  Denote by $\log$ the standard
branch of the logarithm, which has the imaginary part in
$]-\pi,\pi]$. Recall that the space of triples of log-branches on a
    branched oriented tetrahedron is identified with the Riemann surface
    $\widehat{\mc}$. 
\smallskip

For $N=1$, we forget the integral charge $c$, so that
$\mathcal{R}_1$ is defined on flattened $\Ii$-tetrahedra:
\begin{equation}\label{symRmatdil}
\mathcal{R}_1(\Delta,b,w,f) =
\exp\biggl(\frac{*_b}{\pi\sqrt{-1}}\mathcal{R}(w_0;f_0,f_1)\biggr) =
\exp\biggl(\frac{*_b}{\pi\sqrt{-1}}\biggl(-\frac{\pi^2}{6} -
\frac{1}{2} \int_0^{w_0} \biggl( \frac{{\rm l}_0(t)}{1-t} - \frac{{\rm
    l}_1(t)}{t} \biggr) \ dt \biggr) \biggr)
\end{equation}
where ${\rm l}_0(t)=\log(t) + \sqrt{-1}\pi f_0$ and ${\rm
  l}_1(t)=\log((1-t)^{-1}) + \sqrt{-1}\pi f_1$. The map
  $\mathcal{R}:\widehat{\mc}\rightarrow \mc/\pi^2\mz$ is holomorphic,
  and takes values in $\mc/2\pi^2\mz$ on the component with even
  valued flattenings (\cite{N}, Proposition 2.5).  \smallskip

For $N=2m+1>1$ and every complex number $x$ set $x^{1/N} =
\exp(\log(x)/N)$ ($0^{1/N} = 0$ by convention). Put
\begin{equation}\label{funcg}
g(x) := \prod_{j=1}^{N-1}(1 - x\zeta^{-j})^{j/N}.
\end{equation}
The function $g$ is defined over $\mc$, and analytic over the
complement of the rays from $x =\zeta^k$ to infinity, $k=1,\ldots,\
N-1$. Set $h(x) := g(x)/g(1)$ (we have $\vert g(1) \vert = N^{1/2}$).
For any $u',v'\in \C$ satisfying $(u')^N + (v')^N = 1$ and any $ n \in
\mathbb{N}$, let
$$\omega(u',v'\vert n) = \prod_{j=1}^n \frac{v'}{1-u'\zeta^j}$$ with
$\omega(u',v'\vert 0) = 1$ by convention. The functions $\omega$ are
periodic in their integer argument, with period $N$. Write
$[x]=N^{-1}(1-x^N)/(1-x)$. Given a flat/charged $\Ii$-tetrahedron
$(\Delta,b,w,f,c)$, set
\begin{equation}\label{qvar}
w_j'=\exp({\rm l}_{j,N}/N)
\end{equation}
with ${\rm l}_{j,N}$ as in (\ref{qlb}). Put the standard tensor
product basis on $\C^N\otimes \C^N$. The matrix dilogarithm of level
$N>1$ is the tensor valued function of flat/charged $\Ii$ tetrahedra
defined by
\begin{equation}\label{symqmatdil}
\mathcal{R}_N(\Delta,b,w,f,c) =
\bigl((w_0')^{-c_1}(w_1')^{c_0}\bigr)^{\frac{N-1}{2}}\
(\Ll_N)^{*_b}(w_0',(w_1')^{-1})\in {\rm
Aut}(\C^N\otimes\C^N)
\end{equation}
where (recall that $N=2m+1$)
$$\begin{array}{lll} \Ll_N(u',v')_{k,l}^{i,j}& = & h(u')\
\zeta^{kj+(m+1)k^2}\ \omega(u',v'\vert i-k) \ \delta(i + j - l) \\
\bigl( \Ll_N(u',v')^{-1}\bigr)_{k,l}^{i,j}& = & \frac{[u']}{h(u')}\
\zeta^{-il-(m+1)i^2}\ \frac{\delta(k + l - j)}{\omega(u'/\zeta,v'\vert
k-i)}
\end{array}$$
with $\delta$ the Kronecker symbol with period $N$, that is $\delta(n)
= 1$ if $n \equiv 0$ mod($N$), and $\delta(n) = 0$ otherwise. Note
that we use the branching in order to associate an index among $i$,
$j$, $k$ and $l$ to each $2$-face of $\Delta$. The rule is shown in
Figure \ref{CQDidealtensor}.

Up to multiplication by $N$th roots of unity, the map $\mathcal{R}_N$
is holomorphic on $\widehat{\mc}$, because $w_j'= \exp((1/N) ({\rm
  l}_j - *_bc_j)) \exp((f_j - *_bc_j)\pi i)$ and the parity of $f_j$
is unaltered when we move $l_j$ continuously. The ambiguity of
$\mathcal{R}_N$ is a consequence of the jumps of $g$ along the cuts
from the $\zeta^k$ to infinity.
\medskip

\noindent {\bf On matrix dilogarithms via geometric quantization.}  As
explained in Remark \ref{OKCCS}, in this paper we concentrate on the
{\it quantum} matrix dilogarithms (ie.  with $N>1$). These were
derived in \cite{BB2} from the Kashaev's 6j-symbols for the cyclic
representation theory of a Borel quantum subalgebra $\Bb_\zeta$ of
$U_\zeta(sl(2,\mc))$, where $\zeta=\exp(2i\pi/N)$. Let us outline here
very briefly an alternative construction based on {\it geometric
quantization of $\Rr_1$}, thus clarifying their geometric origin (for
details see \cite{B}).
\smallskip

Consider an abstract oriented quadrilateral $Q$, triangulated by two
triangles.  Order the triangles of $Q$, and associate to each a copy
$X_i$, $i=1,2$, of $X=\{(u,v,w)\in(\mc^*)^3\mid uvw=-1\}$, where $u$,
$v$ and $w$ correspond to the corners, ordered cyclically by using the
orientation. Let us regard $\mc$ as the Riemann surface of $\log$,
ie. 
$$\mc=\{(z;p) \in \mc^*\times 2\mz \}/((z+i0;p)\sim (z-i0;p+2),
\forall z\in (-\infty;0)) \ . $$ As in (\ref{clb}), we set ${\rm
l}(z;p)= \log(z)+\sqrt{-1}\pi p $.  Let
$$\widehat{X}=\{((u;p),(v;q),(w;r)) \in \tilde{X}\mid {\rm l}(u;p)+{\rm
l}(v;q)+{\rm l}(w;r)=0\}$$ a subspace of the universal covering
$\tilde{X}$ of $X$, and put
$$\widehat{\mc}(z;p,q)=\{({\rm l}(z;p),{\rm
l}((1-z)^{-1};q))\}=\{((z;p),((1-z)^{-1};q),(1-z^{-1};r))\in
\widehat{X}\} \ . $$ Denote $b_Q=d\log u_1 \wedge d\log v_1 + d\log
u_2 \wedge d\log v_2$ the canonical complex symplectic form on
$X_1\times X_2$, and
$$\theta_Q=(1/2)\left(({\rm l}(u_1;p_1)d\log v_1 - {\rm
  l}(v_1;q_1)d\log u_1) + ({\rm l}(u_2;p_2)d\log v_2 - {\rm
  l}(v_2;q_2)d\log u_2))\right)$$
the symplectic potential for the lift of $b_Q$ to $\widehat{X}_1
\times \widehat{X}_2$. Consider the four punctured sphere $S_Q^2$
obtained by gluing along the boundary in the natural way $Q$ and the
quadrilateral $Q'$ obtained from it by exchanging the diagonal. The
$S_4$-action on the vertices of $S_Q^2$ reorders the copies $X_i$
($i=1,\ldots,4$) attached to the triangles, and induces the usual action
of $PSL(2,\mz)$ on each $\widehat{X}_i\cong \mc^2$, with standard
basis ${\rm l}(u_i;p_i)$ and ${\rm l}(v_i;q_i)$, by
symplectomorphisms. 
\begin{prop}\label{flipform} 
There is a canonical family $\{\tilde{\phi}_Q^n\}_{n\in \mz}:
(\widehat{X}_1 \times \widehat{X}_2,b_Q)\rightarrow (\widehat{X}_3
\times \widehat{X}_4,b_{Q'})$ of analytic symplectomorphisms
equivariant for the $S_4$-action on $S_Q^2$, given by (see
Figure \ref{diagex})
$$ \tilde{\phi}_Q^n(({\rm l}(u_1;p_1),{\rm l}(v_1;q_1)),({\rm
    l}(u_2;p_2),{\rm l}(v_2;q_2))) = (({\rm l}(u_1';p_1'),{\rm
    l}(v_1';q_1')),({\rm l}(u_2';q_2'),{\rm l}(v_2';q_2'))),$$ where we
    identify $\widehat{X}$ with $\mc^2$ (first two coordinates) and
$$\left\lbrace\begin{array}{lll}{\rm l}(u_1';p_1') & = & {\rm
l}(u_1;p_1) -{\rm l}(u_2';p_2') \\{\rm l}(v_1';q_1')& = & {\rm
l}(v_1;q_1) + {\rm l}(u_2;p_2)\end{array}\right.\quad
\left\lbrace\begin{array}{lll} {\rm l}(u_2';p_2') & = & {\rm
l}(u_2;p_2)-{\rm l}((1-u_1v_2)^{-1};n)\\{\rm l}(v_2';q_2')& = & {\rm
l}(v_2;q_2) +{\rm l}(u_1';p_1')+{\rm l}(v_1';q_1')
\end{array}\right. .$$
The maps $\tilde{\phi}_Q^n$ satisfy the pentagon relation, that is,
$\tilde{\phi}_{Q_1}^{n_1}\circ\tilde{\phi}_{Q_2}^{n_2}
\circ\tilde{\phi}_{Q_3}^{n_3}\circ(\tilde{\phi}_{Q_4}^{n_4})^{-1}\circ
(\tilde{\phi}_{Q_5}^{n_5})^{-1}={\rm Id}_{\widehat{X}^3}$, where $Q_i$
has diagonal the $i$th edge exchanged in Figure \ref{figpent}, for the
positive cyclic ordering starting from the top left pentagon. Moreover
$\theta_Q -\tilde{\phi}_n^*\theta_{Q'}= -d\mathcal{R}$, where
$\mathcal{R}:\widehat{\mc}(u_1v_2;m,n)\rightarrow \mc/2\pi\mz$ is the
extended dilogarithm of (\ref{symRmatdil}), and
$\widehat{\mc}(u_1v_2;m,n)=\{({\rm l}(u_1;p_1)+{\rm l}(v_2;q_2),{\rm
l}((1-u_1v_2)^{-1};n))\}$ is attached to the diagonal of $Q$.
\end{prop}
By ``canonical'' we mean that $\{\tilde{\phi}_Q^n\}$ is the unique
satisfying some natural properties related to configuration spaces of
points in $\mc P^1$. In fact, $\widehat{\mc}(z;p,q)$ is isomorphic to
the moduli space of similarity classes of triangles in the complex
plane endowed with lifts to $\mr$ of the angles, or, equivalently, to
the moduli space of isometry classes of hyperbolic ideal tetrahedra
with even valued flattenings. The complex symplectic form $d\log u
\wedge d\log v$ restricts to the {\it real} symplectic form
$w=-d\log(z) \wedge d\log(1-z)$ on $\widehat{\mc}(z;p,q)$. It is the
differential version of the extended complex Dehn invariant
$\widehat{\delta}:\widehat{\mc}(z;p,q) \rightarrow \mc \wedge \mc$,
$(z;p,q)\mapsto {\rm l}(z;p)\wedge_\mz {\rm l}(1-z;q)$ of \cite{N}.

When Im$(z)\ne 0$ the form $w$ is K\"ahler for the usual complex
structure on $\mr^2$. Replace $\widehat{\mc}$ with the ramified
covering $\widehat{\mc}_N$ of $\mc^*$ obtained by taking the quotient
with $N\mz\times N\mz$. Put the forms $Nw$ and $N\theta$ on
$\widehat{\mc}_N$. Standard half-form quantization of
$\widehat{\mc}_N$ produces a $N$-dimensional vector space $\Gamma^N$
of sections of a line bundle over $\widehat{\mc}_N$ (the {\it coherent
  states}), and the maps $\tilde{\Phi}_Q^n:
\Gamma_1^N\otimes\Gamma_2^N\rightarrow \Gamma_3^N\otimes \Gamma_4^N$
induced by pull back via $\{\tilde{\phi}_Q^n\}$, that is such that
$\tilde{\Phi}_Q^n(s)=s\circ \tilde{\phi}_Q^n$ for any $s \in
\Gamma_1^N\otimes\Gamma_2^N$, coincide with the matrix dilogarithms
$\Rr_N$ for some suitable basis. In particular, the pentagon relation for
$\tilde{\phi}_Q^n$ lifts to the five term identities mentionned after
Proposition \ref{fund_func} below.
\begin{figure}[ht]
\begin{center}
\includegraphics[width=6cm]{flipchart.eps}
\caption{\label{diagex} A diagonal exchange.}
\end{center}
\end{figure}
\begin{figure}[ht]
\begin{center}
 \includegraphics[width=7cm]{pentfig.eps}
\caption{\label{figpent} The pentagon relation.}
\end{center}
\end{figure}

\subsection{QHG triangulated pseudomanifolds}\label{QHGPM} We
restrict the discussion to pseudomanifolds for simplicity, but all
what follows makes sense for arbitrary singular $3$-cycles whose non
manifold locus is of codimension $\geq 2$. By a pseudomanifold $Z$,
possibly with non-empty boundary $\partial Z$, we mean a compact
oriented polyhedron with at most a finite set of non manifold
points. The boundary is a pseudo-surface.

A {\it QHG triangulated pseudomanifold} $(Z,\Tt)$ is a pseudomanifold
$Z$ obtained as the quotient of a finite family
$\mathcal{Z}=\{(\Delta^i,b^i,w^i,f^i,c^i)\}$ of flat/charged
$\Ii$-tetrahedra, via a system of orientation reversing simplicial
identifications of pairs of $2$-faces such that the branchings
match. The resulting triangulation $T$ of $Z$ is endowed with a
{\it global branching} $b$, and is possibly singular (multiply
adjacent as well as self adjacent tetrahedra are allowed). The set of
non manifold points of $Z$ is contained in the set of vertices of
$T$. We do not impose for the moment any global constraint on the
moduli, flattening and charges. Hence $Z$ is equipped with a {\it
rough flat/charged $\Ii$-triangulation} $\ \Tt=(T,b,w,f,c)$, where
$w=\{w^i\}$ and so on.

Next we define the {\it QHG triangulated pseudomanifold
isomorphisms}. Fix a QHG triangulated pseudomanifold $(Z,\Tt)$,
and let $\epsilon_T : E(\Zz)\to E(T)$ be the identification map of
edges. We define the {\it total modulus, log-branch} and {\it
charge} of an edge $e$, respectively, by:
\begin{eqnarray}\label{totw} W_\Tt(e)=\prod_{h \in
\epsilon_T^{-1}(e)} w(h)^{*_{b}}\nonumber\\ L_\Tt(e)= \sum_{h\in
\epsilon_T^{-1}(e)} *_{b}\ {\rm l}(h)\\ C_\Tt(e)= \sum_{h\in
\epsilon_T^{-1}(e)} c(h)\ \ \ \nonumber
\end{eqnarray}
where $*_{b}=\pm 1$ according to the $b$-orientation of the
tetrahedron in $Z$ that contains $h$, $w(h)$ is the cross-ratio
modulus at $h$, $l(h)$ the log-branch at $h$, and $c(h)$ the
charge at $h$.
\begin{remark}\label{W_mean} {\rm It is easily seen that $W_\Tt(e)$ is
a cross ratio for the four `extremal' points on
$\partial \overline{\mathbb{H}}^3$ determined by gluing oriented
hyperbolic ideal tetrahedra with moduli $w(h)^{*_{b}}$ along a
common edge (by continuity $W_\Tt(e)=1$ for a degenerate
quadrilateral with
    three distinct vertices).}
\end{remark}
It is well-known that any two arbitrary {\it naked} triangulations $T$,
$T'$ of $Z$ with the same boundary triangulation can be connected
(keeping the boundary triangulation fixed) by a finite sequence of the
{\it local moves} shown in Figure \ref{CQDfigmove1}, the {\it $2
\leftrightarrow 3$ move} (top) and the {\it bubble move} (bottom).

\begin{figure}[ht]
\begin{center}
\includegraphics[width=6cm]{QHFTmove.eps}
\caption{\label{CQDfigmove1} The moves on naked singular
triangulations.}
\end{center}
\end{figure}
For any such a local move $T \leftrightarrow T'$ we have two
triangulations of a same portion of a polyhedron $Q$. Assume that
both $T$ and $T'$ extend to portions $(Q,\Tt)$ and $(Q,\Tt')$ of
QHG triangulated pseudomanifolds. We have to specify the
admissible {\it QHG transits} $\Tt \leftrightarrow \Tt'$. In any
case we require that they are local (that is, the portions
complements remain unchanged), and that the branchings coincide at
every common edge of $T$ and $T'$.

For the $2 \leftrightarrow 3$ move we also require that at every
common edge $e$ as above the total modulus, log-branch and charge
coincide
\begin{equation}\label{preserveq}
W_\Tt(e)=W_{\Tt'}(e),\quad L_\Tt(e)=L_{\Tt'}(e),\quad C_\Tt(e)=C_{\Tt'}(e).
\end{equation}
The same rule restricted to the total modulus and total log-branch
holds also for the bubble move; however, the total charge behaves
in a different way: any bubble transit $\Tt \leftrightarrow \Tt'$
includes a marked edge $e$ common to $T$ and $T'$. Referring to
the bottom of Figure \ref{CQDfigmove1} we require that
$C_\Tt(e)=C_{\Tt'}(e)-2$, while for the other two common edges the
total charges are unchanged.

\begin{remark}\label{forced}{\rm
For every QHG transit supported by a $2 \leftrightarrow 3$ move
(top of Figure \ref{CQDfigmove1}), if $E_0$ denotes the new edge
in $T'$ then \begin{equation}\label{2-3comp} W_{\Tt'}(E_0)=1, \
L_{\Tt'}(E_0)=0, \ C_{\Tt'}(E_0)=2.\end{equation} For every QHG
bubble transit, let $f'$ be the unique new $2$-simplex of $T'$
that contains the marked edge $e$. Denote by $E_1$ and $E_2$ the
other edges of $f'$, and by $E_3$ the further new edge of $T'$.
Then we have} \begin{eqnarray}\label{bubblecomp}
 W_{\Tt'}(E_j)=1, \quad  L_{\Tt'}(E_j)=0, \quad j=1,2,3 \\
C_{\Tt'}(E_1)=C_{\Tt'}(E_2)= 0, \quad C_{\Tt'}(E_3)=2.\nonumber
\end{eqnarray}
\end{remark}

\begin{defi}\label{QHGisom}{\rm A {\it QHG isomorphism} between
    QHG triangulated pseudomanifolds is any finite composition of QHG
    transit configurations and oriented simplicial homeomorphisms
that preserve the whole decoration.}
\end{defi}

\subsection{QHG universal functor}\label{QHGFUNC} For every odd
$N\geq 1$, we associate to every QHG triangulated pseudomanifold
$(Z,\Tt)$ a {\it trace tensor} $\Hh_N(\Tt)$, as follows.  Define
an {\it $N$-state} of $T$ as a function that gives every
$2$-simplex an index, with values in $\{0,\dots,N-1\}$. Every
$N$-state determines an entry for each matrix dilogarithm
$\Rr_N(\Delta, b, w, f,c)$. As two tetrahedra induce opposite
orientations on a common $2$-face, an index is down for the
$\Rr_N$ of one tetrahedron while it is up for the other (see
Figure \ref{CQDidealtensor}). By summing over repeated indices we
get the {\it total contraction} of the tensors $\{\Rr_N(\Delta, b,
w, f,c)\}$, that we denote $\textstyle \prod_{\Delta \subset T}
\Rr_N(\Delta,b,w,f,c)$. Let $v_I$ and $v_\delta$ be the number of
vertices of $T\setminus \partial T$ and $\partial T$,
respectively, that correspond to manifold points. We set
\begin{equation}\label{deftt}
\Hh_N(\Tt) = N^{-(v_\delta/2 + v_I)}\prod_{\Delta \subset T}
\Rr_N(\Delta,b,w,f,c)
\end{equation}
The type of a trace tensor $\Hh_N(\Tt)$ depends on the $b$-signs
of the boundary triangles of $(T,b)$. The matrix dilogarithms
themselves are special instances of trace tensors.

We have
\begin{prop}\label{fund_func}
For every odd $N\geq 1$, up to sign and multiplication by $N$th
roots of unity the trace tensor $\Hh_N(\Tt)$ is invariant up to
QHG isomorphism.
\end{prop}
This result is a restatement of Theorem 2.1 (2) and Lemma 6.7 of
\cite{BB2}, and summarizes the fundamental functional relations
satisfied by the matrix dilogarithms. In particular those
corresponding to $2 \leftrightarrow 3$ QHG transits are usually called
{\it five terms identities}. In Figure \ref{CQDSchaeffer} we show one
instance in dual terms. The normalization factor $ N^{-v_I}$ in
(\ref{deftt}) is due to the bubble move, that changes by $1$ the
number of internal vertices.
\begin{figure}[ht]
\begin{center}
 \includegraphics[width=7cm]{CQDShaeffernew.eps}
\caption{\label{CQDSchaeffer} A $2 \leftrightarrow 3$ QHG move
($x_1=y/x$, $x_2=y(1-x)/x(1-y)$, and $x_3=(1-x)/(1-y)$).}
\end{center}
\end{figure}
\medskip

The notions of ``boundary'' (also allowing only portions of the
standard boundary - see eg. \cite{T}), bordism, and bordism gluing are
well defined for the category of QHG triangulated pseudomanifolds
considered up to QHG isomorphisms . Hence, for every odd $N\geq 1$,
the association of the trace tensor to each bordism defines a
functor. Note that the identities between trace tensors hold up to the
phase ambiguity of Proposition \ref{fund_func}. Also, the
normalization factor $N^{-v_\delta/2}$ in (\ref{deftt}) compensates
the change in the number of internal vertices when gluing along
complete {\it connected} components of the (usual) boundary (this is
easily adapted to more general gluings).

Any bordism is the gluing of flat/charged $\Ii$-tetrahedra,
considered as elementary bordisms between the couples of $2$-faces
having $b$-sign equal to $-1$ or $+1$. The matrix dilogarithms can
be interpreted as the amplitudes of a diagonal exchange that
relates two quadrilateral triangulations (recall the above discussion
about the geometric quantization derivation of matrix dilogarithms.)

\subsection{Towards geometric specializations}\label{ToGeoSpec} The
category of QHG triangulated pseudomanifolds is built {\it ad hoc}
on the functional properties of the matrix dilogarithms. In order
to eventually get specializations with a geometric content, such
as QHFT, it is necessary to refine more and more our {\it rough}
flat/charged $\Ii$-triangulations. In fact we have to impose some
{\it global constraints} that would be preserved by (possibly
refined) QHG isomorphisms.  It is useful to recall at once some of
these refinements, and related matter. We adopt the notations
introduced above.
\begin{defi}\label{non-rough}{\rm A triple $\Tt=(T,b,w)$ is an {\it
      $\Ii$-triangulation} if at each edge $e$ not contained in
    $\partial T$ we have the {\it edge compatibility relation}
    $W_\Tt(e)=1$.  We say that $\Tt=(T,b,w,f)$ is a {\it flattened}
    $\Ii$-triangulation (and $f$ a {\it global  flattening}) if moreover
    $L_\Tt(e)=0$.}
\end{defi}
It is easily seen that the collection of cross-ratio moduli of an
$\Ii$-triangulation $\Tt$ of a pseudomanifold $Z$ defines a PL {\it
pseudo developing map} $d: \tilde{Z} \to \overline{\mh}^3$, unique up to
post composition with the action of $PSL(2,\C)$, and a {\it holonomy
representation} $h: \pi_1(Z) \to PSL(2,\C)$ such that $d(\gamma(y)) =
h(\gamma)(d(y))$ for every $\gamma \in \pi_1(Z)$, $y\in \tilde{Z}$.

Our main tool for producing $\Ii$-triangulations is the following
  {\it idealization procedure}:
\begin{defi}\label{3dim-id}{\rm
Let $\Tt=(T,b,z)$ be a branched triangulation of a pseudomanifold
$Z$ equipped with a $PSL(2,\C)$-valued $1$-cocycle $z$, where the
cocycle relation is $z(e_0)z(e_1)z(e_2)^{-1} = 1$ on a branched
triangle with edges $e_0$, $e_1$ and $e_2$. A $3$-simplex $\Delta$
of $(T,b,z)$ with vertices $x_0,x_1,x_2,x_3$ is {\it idealizable}
if
\begin{equation}\label{vertid}
u_0=0,\ u_1= z(x_0x_1)(0),\ u_2= z(x_0x_2)(0),\ u_3= z(x_0x_3)(0)
\end{equation}
are four distinct points in $\mc \subset \mc\mathbb{P}^1= \partial
\overline{\mh}^3$. The convex hull is a hyperbolic ideal tetrahedron
with ordered vertices, denoted $(u_0,u_1,u_2,u_3)$. Let us give
each edge $e$ of $\Delta$ the cross-ratio modulus $w(e)\in \C
\setminus \{0,1\}$ of the corresponding edge in
$(u_0,u_1,u_2,u_3)$, and put $w=(w_0,w_1,w_2) \in (\C \setminus
\{0,\ 1 \})^3$. We call $(\Delta,b,w)$ the {\it idealization} of
$(\Delta,b,z)$. If all $3$-simplices of $(T,b,z)$ are idealizable
we say that $\Tt_{\Ii}=(T,b,w)$ is the {\it idealization} of
$\Tt$.}
\end{defi}
Note that the idealization includes the choice of a base point
(ie. $0$ in Definition \ref{3dim-id}; we make this choice once for
ever).  The idealization of each single $3$-simplex strongly
depends on this choice, but taking a different base point is
equivalent to conjugating the cocycle $z$. Not all branched
triangulations support idealizable cocycles (for example an ideal
triangulation of a knot in $S^3$ only supports the trivial
constant cocycle). However, when it exists we have:
\begin{lem}\label{idealization}{\rm (\cite{BB1}, section 2.4)}
The idealization $\Tt_\Ii$ is an $\Ii$-triangulation of $Z$.
\end{lem}
The following fact will be important below (a version of it was
implicitely used in \cite{BB1}; see (3), Lemma 2.12 and \S 3.2 in
that paper):
\begin{lem}\label{canonicalflat} The idealization $\Tt_\Ii$ is
canonically flattened by taking for each $3$-simplex the
log-branches (with the notations of Definition \ref{3dim-id}):
$$\begin{array}{l}{\rm
l}_0:=\log(u_2-u_1)+\log(u_3)-\log(u_2)-\log(u_3-u_1)\\ {\rm
l}_1:=\log(u_2)+\log(u_3-u_1)-\log(u_0)-\log(u_3-u_2)+\sqrt{-1}\pi\\{\rm
l}_2:=\log(u_3-u_2)+\log(u_0)-\log(u_3)-\log(u_2-u_1)-\sqrt{-1}\pi.
\end{array}$$
Moreover, replacing at each edge of $\Tt_\Ii$ the standard $\log$ with
any other log determination still makes it flattened.
\end{lem}
\noindent {\it Proof.} These expressions are just (corrected)
signed sums of the standard logs of the edge vectors in the
cross-ratio moduli $w_0=(u_2-u_1)u_3/u_2(u_3-u_1)$,
$w_1=u_2(u_3-u_1)/u_0(u_3-u_2)$ and
$w_2=-(u_3-u_2)u_0/(u_2-u_1)u_3$. They clearly define triples of
log-branches. The idealizations for two distinct branchings with
the same $b$-orientation are related by an element in
$PSL(2,\mc)$, that is, a conformal transformation of
$\mathbb{C}P^1$. We deduce that triples $({\rm l}_0,{\rm l}_1,{\rm
l}_2)$ are invariant under a change of branching, because the
angle formed by any pair of vectors is preserved (for instance
$u_2-u_1$ and $u_2$, or $u_3$ and $u_3-u_1$). Then, the edge
compatibility relations $L_\Tt(e)=0$ follow easily by developing
the tetrahedra around $e$ with branchings such that $e=e_0$,
similarly as for Lemma \ref{idealization}. For the last claim we
note that any $\log$ correction appears for each $3$-simplex in
two distinct ${\rm l}_j$, with opposite signs.\hfill$\Box$
\medskip

\noindent{\bf Distinguished flat/charged $\Ii$-triangulations.}
Given an arbitrary non empty and properly embedded tangle $L$ in
$Z$, we say that $(T,H)$ and $\Tt=(T,H,b,w,f,c)$ are {\it
distinguished} if $H$ is a subcomplex of the $1$-skeleton of $T$
isotopic to $L$ and passing through all the vertices that are
manifold points, and containing no singular vertices (we say that
$H$ is Hamiltonian).
\begin{defi}\label{DISFLATCHARGE}{\rm A distinguished $\Ii$-triangulation
$\Tt=(T,H,b,w,f,c)$ is {\it flat/charged} if $f$ is a global flattening,
and
\begin{equation}\label{condc}C_{\Tt}(e)=\left
\lbrace\begin{array}{l} 0\\ 2\end{array}\right.\begin{array}{l}
{\rm if}\ e \subset H\\ {\rm if}\ e \subset T\setminus(H \cup
\partial T)\end{array}
\end{equation} In such a case $c$ is said a {\it global charge}.}
\end{defi}
The existence of global flattenings of an arbitrary $\Ii$-triangulation with
empty boundary, and of global charges (such that $C_\Tt(e)=2$ at all
edges $e$) on any topological ideal triangulation of an oriented
$3$-manifold whose boundary consists of tori, was proved by Neumann in
\cite{N}, section 9, and \cite{N0}, section 6, respectively. This last
result is easily adapted to the existence of global charges on
distinguished triangulations $(T,H)$ (see \cite{BB1}, Theorem 4.7).
\smallskip

We have to refine the QHG isomorphisms in order to deal with
distinguished flat/charged  $\ \Ii$-triangulations. First we have to
incorporate the Hamiltonian tangles into the bare moves. Any
positive $2 \to 3$ move $T\to T'$ naturally specializes to a move
$(T,H)\to (T',H')$; in fact $H'=H$ is still Hamiltonian. For
positive bubble moves, we assume that an edge $e$ of $H$ lies in
the boundary of the involved $2$-simplex $f$; then $e$ lies in the
boundary of a unique $2$-simplex $f'$ of $T'$ containing the new
vertex of $T'$. We define the Hamiltonian subcomplex $H'$ of $T'$
just by replacing $e$ with the other two edges of $f'$. The
inverse moves are defined in the same way; in particular, for
negative $3\to 2$ moves we require that the edge disappearing in
$T$ belongs to $T' \setminus H'$. The $2 \leftrightarrow 3$ QHG
transit specializes verbatim. For the bubble transit we just impose
that the marked edge $e$ (see section \ref{QHGPM}) coincides with
the above edge of $H$. Thanks to (\ref{2-3comp}) and
(\ref{bubblecomp}) we see that distinguished flat/charged
$\Ii$-triangulations are closed under such refined QHG transits.

\begin{remark}\label{univfc}{\rm The residue mod$(2\sqrt{-1}\pi N)$ of
the classical log-branches of a flattened $\Ii$-triangulation are
equivalently given by $N$th roots $w'(h)$ of the cross-ratio moduli
$w(h)$ such that $w'(e_1)w'(e_2)w'(e_3)=1$ at each $3$-simplex and
$\textstyle \prod_{h \in \epsilon_T^{-1}(e)} w'(h)^{*_{b}}=1$ at each
edge, with the notations of (\ref{totw}).  For level $N$ quantum
log-branches of distinguished flat/charged $\Ii$-triangulations this
is replaced with} $$w'(e_1)w'(e_2)w'(e_3)=\exp(-*_b\sqrt{-1}\pi/N)$$
{\rm and} $$\prod_{h \in \epsilon_T^{-1}(e)} w'(h)^{*_{b}}=\left
\lbrace\begin{array}{l} 1 \\ \exp(-2\sqrt{-1}\pi/N)
\end{array}\right.\begin{array}{l} {\rm if}\ e \subset H\\
{\rm otherwise.}\end{array}$$
\end{remark}

\noindent{\bf Cohomological weights and structural facts.} Given
an $\Ii$-triangulation $\Tt=(T,b,w)$ without boundary, let $T_0$
be the complement of an open cone neighborhood of each
$0$-simplex. This is a disjoint union of triangulated closed
oriented surfaces, the links of the vertices in $T$. To each
flattening $f$ of $\Tt$ a class $\gamma(f) \in H^1(\partial
T_0;\mc)$ is associated as follows.  Represent any non zero
integral $1$-homology class $a$ of $\partial T_0$ by ``normal
paths", that is, a disjoint union of oriented essential simple
closed curves transverse to the triangulation and such that no
component enters and exits from the same face of a $2$-simplex.
Such a curve selects a vertex for each $2$-simplex. The value
$\gamma(f)(a)$ is defined as the signed sum of the log-branches of
the edges ending at the vertices selected by $a$. For each vertex
$v$ the sign is $*_b$ if the path goes in the direction given by
the orientation of $\partial T_0$ as viewed from $v$, and $-*_b$
otherwise. Using the edge compatibility relations of Definition
\ref{non-rough} it is easily checked that $\gamma(f)$ does not
depend on the choice of normal path representative. If the
holonomy of $\Tt$ is trivial or parabolic about each (non
manifold) point this class takes values in $2\mz$ (\cite{N},
Proposition 5.2). Similarly we can define $\gamma_2(f) \in
H^1(T_0;\mz/2\mz)$ by using normal paths in $T_0$ and taking
modulo $2$ sum of the flattenings we meet along the paths. We call
$(\gamma(f),\gamma_2(f))$ the {\it (cohomological) weight} of $f$.
These definitions extend immediately to global charge, replacing
log-branches with charges; then branchings are not needed to
compute the signed sums (ie. put $*_b=1$).

In \cite{N0}, section 4 (see also \cite{N}, section 9), Neumann
defines an integral chain complex $\mathcal{J}$ such that global
charges are defined from cycles at level $3$. The maps $\gamma$
and $\gamma_2$ above are well-defined on the third homology group
$H_3(\mathcal{J})$ and satisfy:
\begin{teo}\label{theoweight}{\rm (\cite{N0}, Theorem 5.1)} The sequence
$$0\rightarrow H_3(\mathcal{J})
\stackrel{(\gamma_2,\gamma)}{\longrightarrow} H^1(T_0;\mz/2\mz)
\oplus H^1(\partial T_0;\mz) \stackrel{i^*-r}{\longrightarrow}
H^1(\partial T_0;\mz/2\mz) \rightarrow 0$$ is exact, where $r:
H^1(\partial T_0;\mz) \rightarrow H^1(\partial T_0;\mz/2\mz)$ is
the coefficient map and $i^*: H^1(T_0;\mz/2\mz) \rightarrow
H^1(\partial T_0;\mz/2\mz)$ is induced by the inclusion $\partial
T_0 \rightarrow T_0$.
\end{teo}
In particular, any pair $(h,k) \in H^1(T_0;\mz/2\mz) \times
H^1(\partial T_0;\mz)$ with $r(k)=i^*(h)$ is a weight for some global
charges. This fact extends to log-branches as follows. Assume that the
holonomy of $\Tt$ restricted to each component of $\partial T_0$ takes
values in a Borel subgroup of $PSL(2,\mc)$, so that it fixes some
point in $\partial \overline{\mathbb{H}}^3$. Let $\gamma'$ be defined as
above, except that for each $3$-simplex log-branches are taken with
flattenings $f_0=f_1=0$. It can be checked that $\gamma'(a)$ is the
logarithm of the derivative of the holonomy of $a$ (a similarity), up
to multiples of $2\sqrt{-1}\pi$. Then, any pair $(h,k) \in
H^1(T_0;\mz/2\mz) \times H^1(\partial T_0;\mc)$ such that
\begin{equation}\label{weightcondition}\left\lbrace
\begin{array}{l}(k-\gamma')/\sqrt{-1}\pi \in H^1(\partial T_0;\mz)\\
r((k-\gamma')/\sqrt{-1}\pi)=i^*(h)\end{array}\right.
\end{equation}
 is a weight for some flattening. For instance, when
$\gamma'(a)\in 2\sqrt{-1}\pi\mz$ for all $a$ the first condition
means that $k$ is integral.

Finally, the structure of the spaces of flattenings and integral
charges is given by Theorem 2.4 in \cite{N0}. For each fixed
weight $(h,k)$ they form an affine space over an integral lattice.
Generators have the following combinatorial realization: for each
$3$-simplex in the star of an edge $e$, add $+1$ to the
flat/charges of one of the two other pairs of opposite edges, and
$-1$ for the other pair, so that the total log-branches or charges
stay equal everywhere. In particular, for flattenings of an
idealization $\Tt_\Ii$ any generator is obtained by adding $+1$ to
the log determination at some edge. Hence any flattening of
$\Tt_\Ii$ inducing the weight of the canonical flattening of Lemma
\ref{canonicalflat} differs from it as described in the statement.

We note that the above refined QHG isomorphisms preserve the weights
(see eg. Lemma 4.12 in \cite{BB1}). The difference in considering
global flattenings or charges with different mod($2$) weights in
$H^1(T_0;\mz/2\mz)$ seems to carry not so essential information (see
Theorem \ref{main} and Remark \ref{scissors}). This contrasts with
boundary weights in $H^1(\partial T_0;\mz)$, which play a key role in
surgery formulae (see \cite{N}, theorems 14.5 and 14.7, and section
\ref{DEHNF} below). We note that a process involving $2$-handle
surgery allows to define explicit isomorphisms between lattices of
flat/charges with different boundary weights (see \cite{N}, section
11, p. 457).
